% Section 3: Related Work
% Target length: 600-900 words
% Status: OUTLINE ONLY - each subsection needs prose + \cite{} keys matching references.bib

\section{Related Work}
\label{sec:related-work}

\subsection{Instructional Design: Elaboration and Cognitive Load}
% TODO: expand.
\begin{itemize}
    \item Elaboration Theory \cite{reigeluth1983} sequences instruction from simple
          ``epitomes'' outward, deliberately layering analogies, synthesis, and
          increasing detail onto core ideas --- consistent with what we observed
          Claude doing for high-importance concepts, but never tied to a graph
          metric.
    \item Cognitive Load Theory \cite{sweller1988} and the worked-example effect
          \cite{renkl2014} motivate why foundational concepts (many downstream
          dependents) warrant more germane-load investment: schema failures there
          propagate.
\end{itemize}

\subsection{Concept Maps and ``What to Teach''}
% TODO: expand; verify full citation details for concept-network-analysis paper
% (only saw it via search snippet + 403 on full text -- confirm before submission).
\begin{itemize}
    \item Node-level concept-network analysis has been proposed as an operational
          measure of ``important knowledge'' for curriculum decisions
          \cite{conceptnetwork2025}.
    \item This line of work stops at \emph{what to include/emphasize}; it does not
          propose a quantitative content-length or word-budget formula. That is the
          gap this paper addresses.
\end{itemize}

\subsection{Prerequisite and Curriculum Graphs}
% TODO: expand.
\begin{itemize}
    \item In-degree and out-degree are standard structural features in curriculum /
          prerequisite graph analysis \cite{capire2025, prereqnetworks2024}.
    \item These works target student modeling and course sequencing, not
          per-concept instructional content generation.
    \item \cite{prereqnetworks2024} notes plain in-degree undercounts ``critical
          introductory'' nodes and proposes a reachability-based alternative ---
          directly motivating our move from in-degree to CIS.
\end{itemize}

\subsection{PageRank-Style Recursive Importance}
% TODO: expand.
\begin{itemize}
    \item PageRank \cite{brin1998} popularized recursive importance: a node is
          important if important nodes point to it.
    \item The same idea transfers cleanly beyond the web graph: TextRank
          \cite{mihalcea2004} for keyword/sentence graphs, Eigenfactor
          \cite{bergstrom2007} for citation-network journal influence.
    \item PageRank / random-walk-with-restart scores have already been used as
          \emph{features} for predicting prerequisite relations in educational
          knowledge graphs \cite{sle2019, cikm2016}, but not as a driver of content
          \emph{volume}.
    \item TotalRank \cite{boldi2005} addresses ranking without an arbitrary damping
          factor; we make a related but distinct observation --- that a DAG
          (as opposed to a general graph with cycles) admits an \emph{exact}
          non-iterative solution by construction (\S\ref{sec:formal-definitions},
          Proposition~1).
\end{itemize}

\subsection{Knowledge Space Theory}
% TODO: expand.
\begin{itemize}
    \item Knowledge Space Theory \cite{doignon1985} formalizes prerequisite
          structure and valid learning paths rigorously but does not address
          content volume per node.
\end{itemize}

\subsection{A Contrasting Data Point}
% TODO: expand; verify exact citation for Wikipedia length/PageRank correlation claim.
\begin{itemize}
    \item Wikipedia article length correlates only weakly with PageRank-based
          importance \cite{wikipagerank2016} --- but Wikipedia length is driven by
          organic editor attention, not top-down instructional design, so this is a
          different generative process rather than a contradiction of our claim.
          Worth stating explicitly to preempt the obvious reviewer objection.
\end{itemize}

\subsection{Positioning}
% TODO: 2-3 sentence explicit gap statement, e.g.:
% "No prior work we found combines (a) a transitive/recursive importance metric
%  computed over a prerequisite DAG, with (b) its use to drive per-concept content
%  \emph{volume} (words + non-text elements) in LLM-generated instructional text."
